%% parts/p02/first-order-resolution-supplement.tex --- retained figure
\section*{Editorial resolution figure}
\addcontentsline{toc}{section}{Editorial resolution figure}
\noindent\textit{Reader note.} The exact inherited chapter supplies the governing
resolution definitions, rules, results, and proofs.  The following exact
historical slice retains only the editorial derivation figure and does not
replace or alter the inherited source.

\begin{figure}[htbp]
\centering
\begin{tikzpicture}[
  font=\small,
  node distance=10mm and 18mm,
  clause/.style={draw,rounded corners,align=center,inner sep=4pt},
  >=Stealth]
  \node[clause] (c1) {\(C_1=\{P(x),Q(x)\}\)};
  \node[clause,right=of c1] (c2) {\(C_2=\{\neg P(a)\}\)};
  \node[clause,right=of c2] (c3) {\(C_3=\{\neg Q(a)\}\)};
  \node[clause,below=of c2] (c4) {\(C_4=\{Q(a)\}\)};
  \node[clause,below=of c4] (empty) {empty clause \(\varnothing\)};
  \draw[->] (c1) -- node[left,align=right,font=\footnotesize]
    {Resolution\\\(\{x\mapsto a\}\)} (c4);
  \draw[->] (c2) -- (c4);
  \draw[->] (c4) -- node[left,font=\footnotesize]{Resolution} (empty);
  \draw[->] (c3) -- (empty);
\end{tikzpicture}
\caption{Resolution-refutation DAG.  Each downward node is
the resolvent of its two incoming parents; the final node is the empty clause.
The first parents are already renamed apart because \(C_2\) is ground, and
their MGU is the displayed finite-support substitution.}
\label{fig:p02-fol-resolution-dag}
\end{figure}
\index{resolution!refutation DAG}
